\section{Conclusions}\label{sec:conclusions}

In this work, we have performed simulation-based inference of tSZ cluster counts with a forward model that
runs the cluster finder on simulated multi-frequency maps, for an SO-like survey with the goal-depth
forecast noise of \citet{SO2019}. We have compared it, on the same simulated skies, with catalogue-level
models, adopting the 21-number summary and the form of the catalogue model of \citet{Zubeldia2025},
together with their prior with narrower flat ranges. Our main results are as
follows.

\begin{enumerate}
\item Map-level simulation-based inference of tSZ cluster counts works at the scale of an SO survey on our
simulated skies. To our knowledge, it is the first inference of cluster counts in which every training
simulation runs the cluster finder on simulated multi-frequency maps. At 40.7 per cent of the sky, the
map-level posterior is within 0.11 widths of the truth, with biases of $-0.11\pm0.08$, $+0.10\pm0.04$,
$-0.05\pm0.06$ and $+0.03\pm0.06$ widths in $\sigma_8$, $\omc$, $\asz$ and $\alsz$, respectively
(Table~\ref{tab:bias}). Its 90 and 50 per cent intervals contain the truth for 0.88--0.91 and 0.46--0.51
of 500 prior skies, respectively (Table~\ref{tab:coverage}). At 10 per cent of the sky, it is
within 0.12 widths of the truth, with a 90 per cent coverage of 0.86--0.90 and slightly too narrow $\alsz$
intervals. These results are robust to the training seed (Appendix~\ref{app:robust}), and we discuss the
offset of $+0.10\pm0.04$ widths in $\omc$ in Section~\ref{sec:discussion:omegac}.
\item Map-level inference is affordable. On one GPU, a survey simulation takes 0.42\,s at 10 per cent
of the sky and 1.6\,s at 40.7 per cent, and a training set of 10\,000 surveys takes 70\,min and 4.6\,h,
respectively. At 40.7 per cent of the sky, the whole comparison, with four posteriors of one to five
networks each, takes about 6 GPU hours and 4 CPU hours (Table~\ref{tab:cost}).
\item A catalogue model calibrated on the maps has to carry three effects of the finder, namely confusion
noise, the search lift and neighbours, as well as the upward tail of its errors
(Section~\ref{sec:finder}). After its masking pass, the finder's noise is 4--8 per cent above that of the
instrument and the CMB, because the maps also contain the tSZ signal of faint haloes. In addition, the
finder's search over positions and sizes raises the $q$ of faint clusters, giving a search lift of
$k=1.37$ in equation~(\ref{eq:qobs_fit}) against the finder's own noise. Neighbours widen the scatter
($w=1.09$) and cause nearly all of the lost matches, which amount to 1.0 per cent of the haloes. Finally,
the finder's errors are not Gaussian but have an upward tail (Section~\ref{sec:finder_tail}).
\item The one-point calibration, made on 20 mocks at the true parameters with the finder's own noise, also
fits the search lift, the scatter and the losses. At the true parameters, it is within about 0.2 widths of
the truth in $\sigma_8$, $\omc$ and $\asz$, with $\alsz$ 0.50 widths high at 10 per cent of the sky and
0.52 at 40.7 per cent. Since the halo-by-halo scale factor is 0.990 against the finder's own noise but
0.921 against the noise of the instrument and the CMB, most of the offset is due to the noise
(Section~\ref{sec:finder_fit}). Across the prior, however, the intervals of this model are too narrow or
off-centre. Its 90 per cent intervals contain the true $\sigma_8$ for 0.65 of the skies at 10 per cent of
the sky and for 0.59 at 40.7 per cent (Table~\ref{tab:coverage}).
\item The noise-level calibration, made at 25 points of the prior through the finder's noise level, gives
$\sigma_8$, $\omc$ and $\asz$ with close to correct coverage, 0.86--0.90 at 90 per cent. The slope $\alsz$,
however, stays 0.62 widths high at 10 per cent of the sky and 0.67 at 40.7 per cent, with 90 per cent
coverage of 0.78 and 0.74, respectively. We attribute this to the upward tail of the finder's errors,
which grows with the tSZ signal on the sky and which a Gaussian scatter cannot carry.
\item As a reference, the uncalibrated catalogue model, with a noise table for the instrument and the CMB
only, as in forecasts, is biased on our map-level skies by up to 8.3 widths at 40.7 per cent of the sky
and 5.3 at 10 per cent (Table~\ref{tab:bias}). At these two sizes it places $\sigma_8$ high by 2.85 and
1.86 widths and $\alsz$ low by 8.30 and 5.28 widths, respectively. At the true parameters, the map-level mocks have 9--25 per
cent fewer clusters above $q=8.4$ than this model predicts (Fig.~\ref{fig:counts}). This bias shows the
size of the finder's effects at goal depth with a flat prior on $\alsz$. We expect it to be smaller at higher noise and with informative priors on the tSZ--mass relation, although we have not tested these
cases (Section~\ref{sec:discussion:scope}).
\end{enumerate}

The finder's response should therefore be forward-modelled in an SO cluster-count analysis, at map level or
through a calibration that follows the parameters and is validated against map-level simulations. Before
the simulator used here is applied to data, however, it needs to be extended beyond its present
simplifications, namely the absence of halo clustering and foregrounds, Gaussian stationary noise and an
idealised confirmation (Section~\ref{sec:discussion:limits}).
